\documentclass[11pt,a4paper]{article}
\newcommand{\gradelabel}{Grade 11}
\newcommand{\leconhdr}{Analyse d'un système automatisé}
% Préambule commun aux plans de cours de robotique (Grades 7 à 12).
% Chaque plan de cours commence par \documentclass[11pt,a4paper]{article}, définit
% \gradelabel et \leconhdr (en-têtes de page), puis % Préambule commun aux plans de cours de robotique (Grades 7 à 12).
% Chaque plan de cours commence par \documentclass[11pt,a4paper]{article}, définit
% \gradelabel et \leconhdr (en-têtes de page), puis % Préambule commun aux plans de cours de robotique (Grades 7 à 12).
% Chaque plan de cours commence par \documentclass[11pt,a4paper]{article}, définit
% \gradelabel et \leconhdr (en-têtes de page), puis \input{../../preambule_robotique.tex}.
\usepackage[T1]{fontenc}
\usepackage[utf8]{inputenc}
\usepackage{lmodern}
\usepackage[french]{babel}
\usepackage[margin=2cm,headheight=14pt]{geometry}
\usepackage[most]{tcolorbox}
\usepackage{xcolor,tabularx,booktabs,array,enumitem,fancyhdr,amssymb,tikz}
\usetikzlibrary{arrows.meta,positioning}

\pagestyle{fancy}
\fancyhf{}
\lhead{\small Robotique -- \gradelabel{} -- Novembre 2026}
\rhead{\small \leconhdr}
\cfoot{\thepage}
\setlist{itemsep=1pt,topsep=2pt,leftmargin=*}
\renewcommand{\arraystretch}{1.35}
\newcolumntype{Y}{>{\raggedright\arraybackslash}X}
\newcommand{\cb}{$\square$\,}
\newcommand{\fl}{$\rightarrow$}
\newcommand{\lig}{\noindent\rule{\linewidth}{0.3pt}\\[5mm]}

\newtcolorbox{etape}[3]{enhanced,breakable,colback=#1!6,colframe=#1!70!black,
  coltitle=white,fonttitle=\bfseries,title={#2\hfill #3},boxrule=0.8pt,arc=2mm,
  before skip=8pt,after skip=8pt}
\newtcolorbox{astuce}[1][Différenciation et besoins particuliers]{breakable,
  colback=yellow!12,colframe=yellow!50!black,fonttitle=\bfseries\small,
  title={#1},boxrule=0.5pt,arc=1mm,before skip=6pt,after skip=6pt}
\newtcolorbox{code}[1][Code]{breakable,colback=black!4,colframe=black!60,
  fonttitle=\bfseries\small,title={#1},boxrule=0.5pt,arc=1mm,fontupper=\ttfamily\small,
  before skip=6pt,after skip=6pt}
% Blocs de code coloré (écrire le code sans accents : comme dans l'éditeur)
\newtcblisting{arduino}[1][Code Arduino]{listing only,breakable,colback=black!4,
  colframe=black!60,fonttitle=\bfseries\small,title={#1},boxrule=0.5pt,arc=1mm,
  before skip=6pt,after skip=6pt,
  listing options={language=C++,basicstyle=\ttfamily\small,columns=fullflexible,
  keepspaces=true,showstringspaces=false,commentstyle=\color{green!40!black}\itshape,
  keywordstyle=\color{blue!70!black}\bfseries}}
\newtcblisting{python}[1][Code Python]{listing only,breakable,colback=black!4,
  colframe=black!60,fonttitle=\bfseries\small,title={#1},boxrule=0.5pt,arc=1mm,
  before skip=6pt,after skip=6pt,
  listing options={language=Python,basicstyle=\ttfamily\small,columns=fullflexible,
  keepspaces=true,showstringspaces=false,commentstyle=\color{green!40!black}\itshape,
  keywordstyle=\color{blue!70!black}\bfseries}}
\newcommand{\entete}[2]{%
  \begin{tcolorbox}[colback=#1!10,colframe=#1!70!black,boxrule=0.8pt,arc=2mm]
  {\Large\bfseries #2}\hfill Nom : \rule{4cm}{0.3pt}\quad Date : \rule{2.2cm}{0.3pt}
  \end{tcolorbox}}

% Bandeau de titre : #1 titre, #2 semaine
\newcommand{\titreblock}[2]{%
\begin{tcolorbox}[colback=teal!10,colframe=teal!70!black,arc=2mm,boxrule=1pt]
\centering
{\LARGE\bfseries #1}\\[2mm]
{\large Robotique -- \gradelabel{} -- Plan de cours de 2 heures}\\[1mm]
Enseignant : Francis Ki \quad|\quad #2
\end{tcolorbox}}

% Tableau du minutage : 6 contenus (accroche, mini-leçon, guidée, autonome, mise en commun, exit ticket)
\newcommand{\minutage}[6]{%
\begin{tabularx}{\linewidth}{>{\bfseries}p{2.2cm}p{2.6cm}Y p{3cm}}
\toprule
Minutes & Étape & Contenu & Organisation\\
\midrule
0 -- 10 & Accroche & #1 & Classe entière\\
10 -- 30 & Mini-leçon & #2 & Classe entière\\
30 -- 60 & Pratique guidée & #3 & Collectif puis binômes\\
60 -- 95 & Pratique autonome & #4 & Individuel sur poste\\
95 -- 110 & Mise en commun & #5 & Classe entière\\
110 -- 120 & Exit ticket & #6 & Individuel sur poste\\
\bottomrule
\end{tabularx}}

% Étayage progressif : #1 aide 1, #2 aide 2, #3 aide 3
\newcommand{\aides}[3]{%
\textbf{Étayage progressif de l'aide.}
\begin{center}
\begin{tabularx}{\linewidth}{>{\bfseries}p{2.4cm}Y p{4.4cm}}
\toprule
Niveau d'aide & Contenu & Qui intervient\\
\midrule
Aide 1 & #1 & Le binôme seul\\
Aide 2 & #2 & Le binôme avec la carte d'aide\\
Aide 3 & #3 & L'enseignant\\
\bottomrule
\end{tabularx}
\end{center}}

% Les trois groupes de la pratique autonome : #1 tâche A, #2 tâche B, #3 tâche C
\newcommand{\groupes}[3]{%
\textbf{Consigne commune\,:} chaque élève ouvre le fichier de son groupe, le renomme avec son prénom et travaille seul sur son poste. L'enseignant passe d'abord auprès du groupe A pendant 10 min.

\medskip
\begin{tabularx}{\linewidth}{>{\bfseries}p{3cm}Y p{2.2cm}}
\toprule
Groupe & Tâche sur poste & Annexe\\
\midrule
A : en dessous du niveau & #1 & Annexe C\\
B : au niveau & #2 & Annexe D\\
C : au-dessus du niveau & #3 & Annexe E\\
\bottomrule
\end{tabularx}}

% Mise en commun : #1 groupe A, #2 groupe B, #3 groupe C, #4 conclusion
\newcommand{\miseencommun}[4]{%
\begin{etape}{teal}{5. Mise en commun}{95 -- 110 min}
\begin{enumerate}
 \item (2 min) Retour au calme, écrans mis en veille.
 \item (9 min) Trois passages courts de 3 min, avec valorisation :
 \begin{itemize}
  \item \textbf{Groupe A} : #1
  \item \textbf{Groupe B} : #2
  \item \textbf{Groupe C} : #3
 \end{itemize}
 \item (4 min) Conclusion : #4
\end{enumerate}
\end{etape}}

% Clôture : #1 question 1, #2 question 2
\newcommand{\cloture}[2]{%
\begin{etape}{gray}{6. Clôture : exit ticket}{110 -- 120 min}
Formulaire en ligne ou papier (annexe G) :
\begin{enumerate}
 \item #1
 \item #2
 \item Auto-évaluation : \emph{Je maîtrise / Je progresse / J'ai besoin d'aide.}
\end{enumerate}
Lire les réponses avant la séance suivante pour ajuster les groupes de besoin.
\end{etape}}

% Grille d'observation : #1 #2 #3 intitulés de colonnes
\newcommand{\grille}[3]{%
\subsection*{Grille d'observation (enseignant)}
\begin{tabularx}{\linewidth}{Y p{2.4cm} p{2.4cm} p{2.4cm} Y}
\toprule
Élève & #1 & #2 & #3 & Groupe suivant\\
\midrule
 & & & & \\ & & & & \\ & & & & \\ & & & & \\ & & & & \\
\bottomrule
\end{tabularx}}

% Signature de l'exit ticket
\newcommand{\autoeval}{Je me sens : \cb Je maîtrise \quad \cb Je progresse \quad \cb J'ai besoin d'aide}
.
\usepackage[T1]{fontenc}
\usepackage[utf8]{inputenc}
\usepackage{lmodern}
\usepackage[french]{babel}
\usepackage[margin=2cm,headheight=14pt]{geometry}
\usepackage[most]{tcolorbox}
\usepackage{xcolor,tabularx,booktabs,array,enumitem,fancyhdr,amssymb,tikz}
\usetikzlibrary{arrows.meta,positioning}

\pagestyle{fancy}
\fancyhf{}
\lhead{\small Robotique -- \gradelabel{} -- Novembre 2026}
\rhead{\small \leconhdr}
\cfoot{\thepage}
\setlist{itemsep=1pt,topsep=2pt,leftmargin=*}
\renewcommand{\arraystretch}{1.35}
\newcolumntype{Y}{>{\raggedright\arraybackslash}X}
\newcommand{\cb}{$\square$\,}
\newcommand{\fl}{$\rightarrow$}
\newcommand{\lig}{\noindent\rule{\linewidth}{0.3pt}\\[5mm]}

\newtcolorbox{etape}[3]{enhanced,breakable,colback=#1!6,colframe=#1!70!black,
  coltitle=white,fonttitle=\bfseries,title={#2\hfill #3},boxrule=0.8pt,arc=2mm,
  before skip=8pt,after skip=8pt}
\newtcolorbox{astuce}[1][Différenciation et besoins particuliers]{breakable,
  colback=yellow!12,colframe=yellow!50!black,fonttitle=\bfseries\small,
  title={#1},boxrule=0.5pt,arc=1mm,before skip=6pt,after skip=6pt}
\newtcolorbox{code}[1][Code]{breakable,colback=black!4,colframe=black!60,
  fonttitle=\bfseries\small,title={#1},boxrule=0.5pt,arc=1mm,fontupper=\ttfamily\small,
  before skip=6pt,after skip=6pt}
% Blocs de code coloré (écrire le code sans accents : comme dans l'éditeur)
\newtcblisting{arduino}[1][Code Arduino]{listing only,breakable,colback=black!4,
  colframe=black!60,fonttitle=\bfseries\small,title={#1},boxrule=0.5pt,arc=1mm,
  before skip=6pt,after skip=6pt,
  listing options={language=C++,basicstyle=\ttfamily\small,columns=fullflexible,
  keepspaces=true,showstringspaces=false,commentstyle=\color{green!40!black}\itshape,
  keywordstyle=\color{blue!70!black}\bfseries}}
\newtcblisting{python}[1][Code Python]{listing only,breakable,colback=black!4,
  colframe=black!60,fonttitle=\bfseries\small,title={#1},boxrule=0.5pt,arc=1mm,
  before skip=6pt,after skip=6pt,
  listing options={language=Python,basicstyle=\ttfamily\small,columns=fullflexible,
  keepspaces=true,showstringspaces=false,commentstyle=\color{green!40!black}\itshape,
  keywordstyle=\color{blue!70!black}\bfseries}}
\newcommand{\entete}[2]{%
  \begin{tcolorbox}[colback=#1!10,colframe=#1!70!black,boxrule=0.8pt,arc=2mm]
  {\Large\bfseries #2}\hfill Nom : \rule{4cm}{0.3pt}\quad Date : \rule{2.2cm}{0.3pt}
  \end{tcolorbox}}

% Bandeau de titre : #1 titre, #2 semaine
\newcommand{\titreblock}[2]{%
\begin{tcolorbox}[colback=teal!10,colframe=teal!70!black,arc=2mm,boxrule=1pt]
\centering
{\LARGE\bfseries #1}\\[2mm]
{\large Robotique -- \gradelabel{} -- Plan de cours de 2 heures}\\[1mm]
Enseignant : Francis Ki \quad|\quad #2
\end{tcolorbox}}

% Tableau du minutage : 6 contenus (accroche, mini-leçon, guidée, autonome, mise en commun, exit ticket)
\newcommand{\minutage}[6]{%
\begin{tabularx}{\linewidth}{>{\bfseries}p{2.2cm}p{2.6cm}Y p{3cm}}
\toprule
Minutes & Étape & Contenu & Organisation\\
\midrule
0 -- 10 & Accroche & #1 & Classe entière\\
10 -- 30 & Mini-leçon & #2 & Classe entière\\
30 -- 60 & Pratique guidée & #3 & Collectif puis binômes\\
60 -- 95 & Pratique autonome & #4 & Individuel sur poste\\
95 -- 110 & Mise en commun & #5 & Classe entière\\
110 -- 120 & Exit ticket & #6 & Individuel sur poste\\
\bottomrule
\end{tabularx}}

% Étayage progressif : #1 aide 1, #2 aide 2, #3 aide 3
\newcommand{\aides}[3]{%
\textbf{Étayage progressif de l'aide.}
\begin{center}
\begin{tabularx}{\linewidth}{>{\bfseries}p{2.4cm}Y p{4.4cm}}
\toprule
Niveau d'aide & Contenu & Qui intervient\\
\midrule
Aide 1 & #1 & Le binôme seul\\
Aide 2 & #2 & Le binôme avec la carte d'aide\\
Aide 3 & #3 & L'enseignant\\
\bottomrule
\end{tabularx}
\end{center}}

% Les trois groupes de la pratique autonome : #1 tâche A, #2 tâche B, #3 tâche C
\newcommand{\groupes}[3]{%
\textbf{Consigne commune\,:} chaque élève ouvre le fichier de son groupe, le renomme avec son prénom et travaille seul sur son poste. L'enseignant passe d'abord auprès du groupe A pendant 10 min.

\medskip
\begin{tabularx}{\linewidth}{>{\bfseries}p{3cm}Y p{2.2cm}}
\toprule
Groupe & Tâche sur poste & Annexe\\
\midrule
A : en dessous du niveau & #1 & Annexe C\\
B : au niveau & #2 & Annexe D\\
C : au-dessus du niveau & #3 & Annexe E\\
\bottomrule
\end{tabularx}}

% Mise en commun : #1 groupe A, #2 groupe B, #3 groupe C, #4 conclusion
\newcommand{\miseencommun}[4]{%
\begin{etape}{teal}{5. Mise en commun}{95 -- 110 min}
\begin{enumerate}
 \item (2 min) Retour au calme, écrans mis en veille.
 \item (9 min) Trois passages courts de 3 min, avec valorisation :
 \begin{itemize}
  \item \textbf{Groupe A} : #1
  \item \textbf{Groupe B} : #2
  \item \textbf{Groupe C} : #3
 \end{itemize}
 \item (4 min) Conclusion : #4
\end{enumerate}
\end{etape}}

% Clôture : #1 question 1, #2 question 2
\newcommand{\cloture}[2]{%
\begin{etape}{gray}{6. Clôture : exit ticket}{110 -- 120 min}
Formulaire en ligne ou papier (annexe G) :
\begin{enumerate}
 \item #1
 \item #2
 \item Auto-évaluation : \emph{Je maîtrise / Je progresse / J'ai besoin d'aide.}
\end{enumerate}
Lire les réponses avant la séance suivante pour ajuster les groupes de besoin.
\end{etape}}

% Grille d'observation : #1 #2 #3 intitulés de colonnes
\newcommand{\grille}[3]{%
\subsection*{Grille d'observation (enseignant)}
\begin{tabularx}{\linewidth}{Y p{2.4cm} p{2.4cm} p{2.4cm} Y}
\toprule
Élève & #1 & #2 & #3 & Groupe suivant\\
\midrule
 & & & & \\ & & & & \\ & & & & \\ & & & & \\ & & & & \\
\bottomrule
\end{tabularx}}

% Signature de l'exit ticket
\newcommand{\autoeval}{Je me sens : \cb Je maîtrise \quad \cb Je progresse \quad \cb J'ai besoin d'aide}
.
\usepackage[T1]{fontenc}
\usepackage[utf8]{inputenc}
\usepackage{lmodern}
\usepackage[french]{babel}
\usepackage[margin=2cm,headheight=14pt]{geometry}
\usepackage[most]{tcolorbox}
\usepackage{xcolor,tabularx,booktabs,array,enumitem,fancyhdr,amssymb,tikz}
\usetikzlibrary{arrows.meta,positioning}

\pagestyle{fancy}
\fancyhf{}
\lhead{\small Robotique -- \gradelabel{} -- Novembre 2026}
\rhead{\small \leconhdr}
\cfoot{\thepage}
\setlist{itemsep=1pt,topsep=2pt,leftmargin=*}
\renewcommand{\arraystretch}{1.35}
\newcolumntype{Y}{>{\raggedright\arraybackslash}X}
\newcommand{\cb}{$\square$\,}
\newcommand{\fl}{$\rightarrow$}
\newcommand{\lig}{\noindent\rule{\linewidth}{0.3pt}\\[5mm]}

\newtcolorbox{etape}[3]{enhanced,breakable,colback=#1!6,colframe=#1!70!black,
  coltitle=white,fonttitle=\bfseries,title={#2\hfill #3},boxrule=0.8pt,arc=2mm,
  before skip=8pt,after skip=8pt}
\newtcolorbox{astuce}[1][Différenciation et besoins particuliers]{breakable,
  colback=yellow!12,colframe=yellow!50!black,fonttitle=\bfseries\small,
  title={#1},boxrule=0.5pt,arc=1mm,before skip=6pt,after skip=6pt}
\newtcolorbox{code}[1][Code]{breakable,colback=black!4,colframe=black!60,
  fonttitle=\bfseries\small,title={#1},boxrule=0.5pt,arc=1mm,fontupper=\ttfamily\small,
  before skip=6pt,after skip=6pt}
% Blocs de code coloré (écrire le code sans accents : comme dans l'éditeur)
\newtcblisting{arduino}[1][Code Arduino]{listing only,breakable,colback=black!4,
  colframe=black!60,fonttitle=\bfseries\small,title={#1},boxrule=0.5pt,arc=1mm,
  before skip=6pt,after skip=6pt,
  listing options={language=C++,basicstyle=\ttfamily\small,columns=fullflexible,
  keepspaces=true,showstringspaces=false,commentstyle=\color{green!40!black}\itshape,
  keywordstyle=\color{blue!70!black}\bfseries}}
\newtcblisting{python}[1][Code Python]{listing only,breakable,colback=black!4,
  colframe=black!60,fonttitle=\bfseries\small,title={#1},boxrule=0.5pt,arc=1mm,
  before skip=6pt,after skip=6pt,
  listing options={language=Python,basicstyle=\ttfamily\small,columns=fullflexible,
  keepspaces=true,showstringspaces=false,commentstyle=\color{green!40!black}\itshape,
  keywordstyle=\color{blue!70!black}\bfseries}}
\newcommand{\entete}[2]{%
  \begin{tcolorbox}[colback=#1!10,colframe=#1!70!black,boxrule=0.8pt,arc=2mm]
  {\Large\bfseries #2}\hfill Nom : \rule{4cm}{0.3pt}\quad Date : \rule{2.2cm}{0.3pt}
  \end{tcolorbox}}

% Bandeau de titre : #1 titre, #2 semaine
\newcommand{\titreblock}[2]{%
\begin{tcolorbox}[colback=teal!10,colframe=teal!70!black,arc=2mm,boxrule=1pt]
\centering
{\LARGE\bfseries #1}\\[2mm]
{\large Robotique -- \gradelabel{} -- Plan de cours de 2 heures}\\[1mm]
Enseignant : Francis Ki \quad|\quad #2
\end{tcolorbox}}

% Tableau du minutage : 6 contenus (accroche, mini-leçon, guidée, autonome, mise en commun, exit ticket)
\newcommand{\minutage}[6]{%
\begin{tabularx}{\linewidth}{>{\bfseries}p{2.2cm}p{2.6cm}Y p{3cm}}
\toprule
Minutes & Étape & Contenu & Organisation\\
\midrule
0 -- 10 & Accroche & #1 & Classe entière\\
10 -- 30 & Mini-leçon & #2 & Classe entière\\
30 -- 60 & Pratique guidée & #3 & Collectif puis binômes\\
60 -- 95 & Pratique autonome & #4 & Individuel sur poste\\
95 -- 110 & Mise en commun & #5 & Classe entière\\
110 -- 120 & Exit ticket & #6 & Individuel sur poste\\
\bottomrule
\end{tabularx}}

% Étayage progressif : #1 aide 1, #2 aide 2, #3 aide 3
\newcommand{\aides}[3]{%
\textbf{Étayage progressif de l'aide.}
\begin{center}
\begin{tabularx}{\linewidth}{>{\bfseries}p{2.4cm}Y p{4.4cm}}
\toprule
Niveau d'aide & Contenu & Qui intervient\\
\midrule
Aide 1 & #1 & Le binôme seul\\
Aide 2 & #2 & Le binôme avec la carte d'aide\\
Aide 3 & #3 & L'enseignant\\
\bottomrule
\end{tabularx}
\end{center}}

% Les trois groupes de la pratique autonome : #1 tâche A, #2 tâche B, #3 tâche C
\newcommand{\groupes}[3]{%
\textbf{Consigne commune\,:} chaque élève ouvre le fichier de son groupe, le renomme avec son prénom et travaille seul sur son poste. L'enseignant passe d'abord auprès du groupe A pendant 10 min.

\medskip
\begin{tabularx}{\linewidth}{>{\bfseries}p{3cm}Y p{2.2cm}}
\toprule
Groupe & Tâche sur poste & Annexe\\
\midrule
A : en dessous du niveau & #1 & Annexe C\\
B : au niveau & #2 & Annexe D\\
C : au-dessus du niveau & #3 & Annexe E\\
\bottomrule
\end{tabularx}}

% Mise en commun : #1 groupe A, #2 groupe B, #3 groupe C, #4 conclusion
\newcommand{\miseencommun}[4]{%
\begin{etape}{teal}{5. Mise en commun}{95 -- 110 min}
\begin{enumerate}
 \item (2 min) Retour au calme, écrans mis en veille.
 \item (9 min) Trois passages courts de 3 min, avec valorisation :
 \begin{itemize}
  \item \textbf{Groupe A} : #1
  \item \textbf{Groupe B} : #2
  \item \textbf{Groupe C} : #3
 \end{itemize}
 \item (4 min) Conclusion : #4
\end{enumerate}
\end{etape}}

% Clôture : #1 question 1, #2 question 2
\newcommand{\cloture}[2]{%
\begin{etape}{gray}{6. Clôture : exit ticket}{110 -- 120 min}
Formulaire en ligne ou papier (annexe G) :
\begin{enumerate}
 \item #1
 \item #2
 \item Auto-évaluation : \emph{Je maîtrise / Je progresse / J'ai besoin d'aide.}
\end{enumerate}
Lire les réponses avant la séance suivante pour ajuster les groupes de besoin.
\end{etape}}

% Grille d'observation : #1 #2 #3 intitulés de colonnes
\newcommand{\grille}[3]{%
\subsection*{Grille d'observation (enseignant)}
\begin{tabularx}{\linewidth}{Y p{2.4cm} p{2.4cm} p{2.4cm} Y}
\toprule
Élève & #1 & #2 & #3 & Groupe suivant\\
\midrule
 & & & & \\ & & & & \\ & & & & \\ & & & & \\ & & & & \\
\bottomrule
\end{tabularx}}

% Signature de l'exit ticket
\newcommand{\autoeval}{Je me sens : \cb Je maîtrise \quad \cb Je progresse \quad \cb J'ai besoin d'aide}


\begin{document}
\titreblock{Analyse d'un système automatisé}{Semaine du 2 au 6 novembre 2026}

%==================================================================
\section{Informations générales}
\begin{tabularx}{\linewidth}{>{\bfseries}p{3.6cm}Y}
\toprule
Titre & Analyse d'un système automatisé\\
Classe / effectif & Grade 11, 11 à 20 élèves, niveaux hétérogènes, élèves à besoins particuliers\\
Durée & 2 heures (120 min)\\
Lieu & Salle informatique, un poste par élève\\
Matériel et logiciels &
\begin{itemize}
 \item Un poste par élève avec navigateur web.
 \item Un outil de schéma en ligne : \textbf{diagrams.net} (draw.io, gratuit, sans compte) ou Google Slides.
 \item Vidéoprojecteur ou tableau interactif.
 \item Courtes vidéos ou photos de systèmes automatisés (barrière de parking, porte de supermarché, distributeur de gel, portail) à préparer à l'avance.
 \item Dossier partagé avec les fichiers de départ des trois niveaux.
\end{itemize}\\
Pré-requis &
Systèmes automatisés, chaîne d'information et chaîne d'énergie vus en G8 (révision rapide dans l'accroche). Cours de septembre sur l'IA : utilisation responsable.\\
\bottomrule
\end{tabularx}

\begin{astuce}[Préparation avant le cours (environ 45 min)]
\begin{enumerate}
 \item Sélectionner trois systèmes (un par niveau) et préparer pour chacun une photo légendée ou une vidéo de 30 secondes.
 \item Préparer un diagramme modèle dans draw.io avec les blocs vides (capteur, traitement, actionneur).
 \item Préparer les fichiers des annexes C à E et le formulaire de sortie (annexe G).
 \item Rappeler aux élèves qu'une réponse générée par IA doit être vérifiée et reformulée : l'analyse du système doit rester personnelle.
\end{enumerate}
\end{astuce}

%==================================================================
\section{Objectifs d'apprentissage}
\begin{tcolorbox}[colback=gray!8,colframe=gray!60!black,arc=1mm,boxrule=0.5pt]
\textbf{Objectifs du programme}
\begin{itemize}
 \item Identifier les fonctions d'un système.
 \item Décrire les flux d'information.
\end{itemize}
\textbf{À la fin de la séance, l'élève est capable\,:}
\begin{itemize}
 \item d'énoncer la fonction globale d'un système (ce qu'il fait, sur quoi, pour qui)\,;
 \item de distinguer capteurs, unité de traitement, préactionneurs et actionneurs\,;
 \item de représenter sur un schéma les flux d'information et d'énergie d'un système automatisé.
\end{itemize}
\textbf{Compétences numériques\,:} réaliser un schéma de blocs, nommer et organiser ses fichiers, sourcer une illustration.
\end{tcolorbox}

%==================================================================
\section{Déroulé minuté (120 min)}
\minutage{Que se passe-t-il quand la voiture arrive devant la barrière ?}{Fonction globale, capteurs, traitement, actionneurs}{Le portail automatique en collectif, puis le distributeur de gel en binômes}{Trois niveaux de besoin}{Valoriser chaque groupe}{Trois questions sur poste}

%------------------------------------------------------------------
\begin{etape}{orange}{1. Accroche : la barrière de parking}{0 -- 10 min}
\begin{enumerate}
 \item (5 min) Projeter la vidéo d'une barrière de parking qui s'ouvre quand une voiture arrive. \emph{Que perçoit la barrière ? Qui «~décide~» ? Qu'est-ce qui bouge ?}
 \item (5 min) Recueillir les mots au tableau en trois colonnes : \emph{ce qu'elle perçoit}, \emph{ce qui décide}, \emph{ce qui agit}. Question-problème : \textbf{«~Comment décrire un système automatisé de façon organisée ?~»}
\end{enumerate}
\end{etape}

%------------------------------------------------------------------
\begin{etape}{blue}{2. Mini-leçon : la structure d'un système automatisé}{10 -- 30 min}
\textbf{Support\,:} annexe A. Quatre temps de 5 min.

\textbf{Temps 1 : la fonction globale.} Une phrase : \emph{le système fait quoi, sur quoi, pour qui}. Exemple : la barrière laisse passer ou arrête les véhicules (sur les véhicules) pour les usagers du parking.

\textbf{Temps 2 : la partie commande et la partie opérative.}
\begin{center}
\begin{tikzpicture}[box/.style={draw,rounded corners,thick,minimum height=1.1cm,text width=2.7cm,align=center,font=\small},>={Stealth[length=3mm]},node distance=0.8cm]
\node[box,fill=orange!20] (a) {\textbf{Capteurs}\\ détectent};
\node[box,fill=blue!15,right=of a] (b) {\textbf{Traitement}\\ décide};
\node[box,fill=green!15,right=of b] (c) {\textbf{Préactionneur}\\ distribue};
\node[box,fill=violet!15,right=of c] (d) {\textbf{Actionneur}\\ agit};
\draw[->,thick] (a)--(b); \draw[->,thick] (b)--(c); \draw[->,thick] (c)--(d);
\node[above=0.1cm of b,font=\scriptsize\itshape] {chaîne d'information};
\node[above=0.1cm of c,xshift=1.7cm,font=\scriptsize\itshape] {chaîne d'énergie};
\end{tikzpicture}
\end{center}

\textbf{Temps 3 : le vocabulaire.}
\begin{center}
\begin{tabularx}{\linewidth}{>{\bfseries}p{3.4cm}Y Y}
\toprule
Élément & Rôle & Exemple (barrière)\\
\midrule
Capteur & Transforme une grandeur physique en information & Boucle magnétique ou cellule\\
Unité de traitement & Compare, décide, envoie des ordres & Microcontrôleur ou automate\\
Préactionneur & Distribue l'énergie à l'actionneur & Contacteur, relais\\
Actionneur & Convertit l'énergie en action & Moteur de la barrière\\
Interface homme-machine & Permet à l'humain de commander ou de voir & Télécommande, bouton\\
\bottomrule
\end{tabularx}
\end{center}

\textbf{Temps 4 : les flux.} Le flux d'\emph{information} circule du capteur vers le traitement puis vers les ordres. Le flux d'\emph{énergie} alimente l'actionneur. Savoir les distinguer permet de comprendre pourquoi un petit signal peut commander un gros moteur.

\begin{astuce}
Utiliser une couleur par rôle (capteur, traitement, actionneur) dans tous les schémas. Fournir un schéma pré-rempli à trous pour les élèves qui ont besoin d'un appui visuel fort. Autoriser une description orale enregistrée.
\end{astuce}
\end{etape}

%------------------------------------------------------------------
\begin{etape}{green}{3. Pratique guidée}{30 -- 60 min}
\textbf{Fiche\,:} annexe B. Les élèves travaillent en binômes devant leur poste.

\textbf{3a. Le portail automatique en collectif (12 min).} Compléter ensemble le tableau d'analyse (fonction globale, capteurs, traitement, préactionneur, actionneur, interface) à partir de la photo ou de la vidéo, puis tracer le schéma dans draw.io en projetant le fichier de l'enseignant.

\textbf{3b. Le distributeur de gel automatique en binômes (18 min).} Réaliser le même travail sur un système neuf : tableau et schéma de blocs. Alternance des rôles (rédacteur, dessinateur).

\aides{Le tableau du portail rempli à l'écran et le schéma modèle}{Questions guides : \emph{Qu'est-ce qui détecte ? Qu'est-ce qui décide ? Qu'est-ce qui agit ? D'où vient l'énergie ?}}{L'enseignant remplit une ligne du tableau avec le binôme, qui poursuit}

\begin{astuce}
Faire commencer toujours par la fonction globale en une phrase. Accepter un croquis à la main photographié si le dessin sur poste ralentit l'élève.
\end{astuce}
\end{etape}

%------------------------------------------------------------------
\begin{etape}{violet}{4. Pratique autonome différenciée}{60 -- 95 min}
\groupes{Barrière de parking : tableau d'analyse à compléter avec une banque de mots (boucle magnétique, microcontrôleur, relais, moteur, bouton, tickets) et schéma de blocs à légender avec des étiquettes mobiles.}{Analyse complète de la porte automatique d'un supermarché : fonction globale, tableau, schéma de blocs (flux d'information et d'énergie) et trois scénarios de fonctionnement (une personne approche, personne ne passe, personne dans l'embrasure).}{Rétro-conception d'une serre automatisée (température, humidité, arrosage) : fonction globale, schéma de blocs, justification du choix de chaque capteur et proposition d'une amélioration (par exemple alerte par message) avec un schéma modifié.}

\medskip
\textbf{Plan B si un élève termine tôt\,:} analyser un système de la maison (lave-linge, réfrigérateur) avec le même tableau.
\end{etape}

%------------------------------------------------------------------
\miseencommun{un élève présente son tableau complété et lit la fonction globale. Valoriser la précision du vocabulaire.}{un binôme projette son schéma et suit le flux d'information avec le doigt. Valoriser la clarté du schéma.}{un élève présente son analyse de la serre et son amélioration. Valoriser le raisonnement technique.}{\textit{tout système automatisé se décrit par sa fonction globale, une chaîne d'information (capter, traiter, ordonner) et une chaîne d'énergie.}}

\cloture{Donne la fonction globale d'une porte automatique en une phrase.}{Quelle est la différence entre un capteur et un actionneur ? Donne un exemple de chaque.}

%==================================================================
% ANNEXES
%==================================================================
\newpage
\section*{Annexe A : Support résumé de la mini-leçon}
\entete{blue}{Analyse d'un système automatisé}

\begin{tcolorbox}[colback=yellow!10,colframe=yellow!50!black]
\textbf{Fonction globale} : le système fait \emph{quoi}, \emph{sur quoi}, \emph{pour qui}.
\end{tcolorbox}

\begin{center}
\begin{tikzpicture}[box/.style={draw,rounded corners,thick,minimum height=1.1cm,text width=2.7cm,align=center,font=\small},>={Stealth[length=3mm]},node distance=0.8cm]
\node[box,fill=orange!20] (a) {Capteurs};
\node[box,fill=blue!15,right=of a] (b) {Traitement};
\node[box,fill=green!15,right=of b] (c) {Préactionneur};
\node[box,fill=violet!15,right=of c] (d) {Actionneur};
\draw[->,thick] (a)--(b); \draw[->,thick] (b)--(c); \draw[->,thick] (c)--(d);
\end{tikzpicture}
\end{center}

\begin{tabularx}{\linewidth}{>{\bfseries}p{3.4cm}Y}
\toprule
Élément & Rôle\\
\midrule
Capteur & Mesure ou détecte, produit une information\\
Unité de traitement & Décide selon un programme\\
Préactionneur & Distribue l'énergie (relais, contacteur)\\
Actionneur & Agit (moteur, lampe, vérin)\\
Interface homme-machine & Commande ou affichage\\
\bottomrule
\end{tabularx}

%==================================================================
\newpage
\section*{Annexe B : Fiche de pratique guidée}
\entete{green}{Pratique guidée : analyser un système}

\subsection*{3a. Le portail automatique (collectif)}
Fonction globale : \dotfill

\begin{tabularx}{\linewidth}{>{\bfseries}p{3.4cm}Y Y}
\toprule
Élément & Dans le portail & Rôle\\
\midrule
Capteur(s) & & \\
Traitement & & \\
Préactionneur & & \\
Actionneur & & \\
Interface homme-machine & & \\
\bottomrule
\end{tabularx}

\subsection*{3b. Le distributeur de gel (binôme)}
Fonction globale : \dotfill

\begin{tabularx}{\linewidth}{>{\bfseries}p{3.4cm}Y Y}
\toprule
Élément & Dans le distributeur & Flux (information ou énergie) \\
\midrule
Capteur(s) & & \\
Traitement & & \\
Actionneur & & \\
Alimentation & & \\
\bottomrule
\end{tabularx}

\begin{tcolorbox}[colback=yellow!10,colframe=yellow!50!black,title={Cartes d'aide à utiliser dans l'ordre},fonttitle=\bfseries\small]
\textbf{Aide 1 :} regarde le tableau du portail et le schéma de l'annexe A.\\
\textbf{Aide 2 :} Qu'est-ce qui détecte ? Qu'est-ce qui décide ? Qu'est-ce qui agit ? D'où vient l'énergie ?\\
\textbf{Aide 3 :} demande à l'enseignant de remplir une ligne avec toi.
\end{tcolorbox}

%==================================================================
\newpage
\section*{Annexe C : Niveau A (en dessous du niveau)}
\entete{orange}{Niveau A : la barrière de parking}

\textbf{Banque de mots :} boucle magnétique, microcontrôleur, relais, moteur, bouton, ticket.

\begin{tabularx}{\linewidth}{>{\bfseries}p{3.4cm}Y}
\toprule
Rôle & Élément de la barrière\\
\midrule
Détecte la voiture & \\
Décide d'ouvrir & \\
Distribue l'énergie & \\
Fait lever la barrière & \\
Commande par l'usager & \\
\bottomrule
\end{tabularx}

\subsection*{Légende le schéma (glisse les étiquettes)}
Capteur \fl{} \dotfill{} \fl{} \dotfill{} \fl{} \dotfill

\subsection*{Vrai ou faux ?}
1. Le moteur est un capteur. \cb Vrai \cb Faux\\
2. Le microcontrôleur décide. \cb Vrai \cb Faux

%==================================================================
\newpage
\section*{Annexe D : Niveau B (au niveau)}
\entete{blue}{Niveau B : la porte du supermarché}

\subsection*{1. Fonction globale}
\lig

\subsection*{2. Tableau d'analyse}
\begin{tabularx}{\linewidth}{>{\bfseries}p{3.4cm}Y Y}
\toprule
Élément & Dans la porte & Grandeur ou énergie concernée\\
\midrule
Capteur(s) & & \\
Traitement & & \\
Préactionneur & & \\
Actionneur & & \\
\bottomrule
\end{tabularx}

\subsection*{3. Trois scénarios}
\begin{tabularx}{\linewidth}{>{\bfseries}p{3.4cm}Y Y}
\toprule
Scénario & Ce que détecte le capteur & Ce que fait la porte\\
\midrule
Une personne approche & & \\
Personne ne passe & & \\
Personne dans l'embrasure & & \\
\bottomrule
\end{tabularx}

Dessine ensuite le schéma de blocs sur une diapositive ou dans draw.io.

%==================================================================
\newpage
\section*{Annexe E : Niveau C (au-dessus du niveau)}
\entete{violet}{Niveau C : rétro-conception d'une serre automatisée}

\begin{tabularx}{\linewidth}{>{\bfseries}p{3.4cm}Y Y}
\toprule
Fonction & Capteur ou actionneur choisi & Pourquoi ce choix ?\\
\midrule
Mesurer la température & & \\
Mesurer l'humidité du sol & & \\
Arroser & & \\
Aérer & & \\
\bottomrule
\end{tabularx}

\subsection*{Amélioration}
Propose une amélioration (par exemple une alerte par message quand la température est trop haute) et indique ce qui change dans le schéma.\\
\lig \lig

\textbf{Défi bonus :} identifie un risque de panne (capteur défectueux) et propose une solution de sécurité.

%==================================================================
\newpage
\section*{Annexe F : Corrections et attendus (enseignant)}

\subsection*{Annexe B}
Portail : fonction globale : ouvrir ou fermer l'accès à un lieu pour les personnes ou véhicules autorisés. Capteurs : cellule photoélectrique, récepteur de télécommande. Traitement : carte de commande. Préactionneur : relais ou variateur. Actionneur : moteur. IHM : télécommande. Distributeur de gel : capteur de présence infrarouge, microcontrôleur, pompe (actionneur), alimentation par piles.

\subsection*{Annexe C}
Boucle magnétique ; microcontrôleur ; relais ; moteur ; bouton ou ticket. Schéma : capteur \fl{} microcontrôleur \fl{} relais \fl{} moteur. Faux, Vrai.

\subsection*{Annexe D}
Fonction globale : permettre ou interdire le passage des personnes à l'entrée du magasin. Capteurs : radar ou cellule infrarouge. Traitement : carte de commande. Actionneur : moteur et rail. Scénarios : approche, ouverture ; personne : fermeture après temporisation ; embrasure : porte maintenue ouverte (sécurité).

\subsection*{Annexe E}
Température : sonde de température ; humidité du sol : capteur d'humidité ; arrosage : électrovanne ; aération : moteur de fenêtre ou ventilateur. Amélioration : module Wi-Fi qui envoie un message. Risque : capteur en panne, d'où un second capteur ou une valeur par défaut sécurisée.\\
\textbf{Critères de réussite :} vocabulaire correct, flux d'information et d'énergie distingués, amélioration cohérente.

%==================================================================
\newpage
\section*{Annexe G : Exit ticket et grille d'observation}
\entete{gray}{Exit ticket : analyse d'un système}

\begin{enumerate}
 \item Donne la fonction globale d'une porte automatique en une phrase.\\ \lig
 \item Quelle est la différence entre un capteur et un actionneur ? Donne un exemple de chaque.\\ \lig
 \item \autoeval
\end{enumerate}
\textbf{Attendus :} 1. Permettre ou interdire le passage des personnes. 2. Le capteur détecte une grandeur et produit une information (par exemple un radar) ; l'actionneur convertit l'énergie en action (par exemple un moteur).

\bigskip
\grille{Fonction globale}{Éléments classés}{Flux représentés}

\end{document}
